\section*{Appendix I: A 45-Day Beachhead Selection Sprint}

\begin{itemize}
    \item \textbf{Days 1 to 7: Map.} Identify 6 to 12 candidate buyer situations. Apply the 6 hard feasibility gates. Rank 3 finalists by urgency, access, delivery feasibility, and evidence quality. List decisive research uncertainties.
    \item \textbf{Days 8 to 21: Reconstruct real decisions.} Interview recent buyers, nonbuyers, and users of substitutes in each finalist group. Reconstruct the actual chronology, money or time spent, and decision authority. Observe workflows where permitted. Remove candidates with fatal barriers.
    \item \textbf{Days 22 to 28: Choose an experiment.} Select the primary candidate. Write a beachhead definition, anti-customer profile, alternative map, concise offer, sample evidence, delivery plan, and downside budget. Predeclare stop and repeat thresholds.
    \item \textbf{Days 29 to 38: Make an actual offer.} Invite qualified buyers through permitted, focused channels. Present a complete product or credible sample alongside a real price. Record decisions and objections. Accept payment only when fulfillment is feasible and clearly contracted.
    \item \textbf{Days 39 to 45: Deliver or diagnose.} Complete the bounded offer where timing permits; otherwise measure the strongest observable commitment and deliver immediately afterward. Record costs, outcomes, and remaining uncertainties. Decide whether to repeat, narrow further, adjust the wedge, or pivot.
\end{itemize}

This is a planning cadence, not a scientific law. A long procurement cycle or multi-month education program may require substantially more time to obtain decision-quality evidence.

\section*{Appendix J: Research Foundations and Further Reading}

These sources support specific concepts; the chapter's decision tables, examples, scorecards, and Bayes recommendations are authorial syntheses and hypotheses. Historical examples are illustrative and subject to survivorship bias. Do not treat celebrated origin stories as causal proof that copying an entry tactic guarantees success.

\begin{enumerate}
    \item[\textbf{[R1]}] \textbf{Geoffrey A. Moore, \emph{Crossing the Chasm} (1991; revised editions).} The beachhead, differences between technology adoption groups, the whole-product model, and bowling-alley expansion. Particularly relevant to technology products moving beyond early adopters.
    \item[\textbf{[R2]}] \textbf{Bill Aulet, \emph{Disciplined Entrepreneurship: 24 Steps to a Successful Startup}, 2nd edition (2024).} Systematically separates market segmentation, beachhead selection, end-user profiling, bottom-up sizing, decision-making units, and subsequent markets.
    \item[\textbf{[R3]}] \textbf{Clayton M. Christensen, Michael E. Raynor, and Rory McDonald (2015), ``What Is Disruptive Innovation?'' \emph{Harvard Business Review}.} Explains low-end and new-market footholds and cautions against labeling every successful entrant a disruptor.
    \item[\textbf{[R4]}] \textbf{Paul Graham (2013), ``Do Things That Don't Scale.''} Founder guidance on manual acquisition, serving the earliest customers with extraordinary care, and concentrating the initial market.
    \item[\textbf{[R5]}] \textbf{Paul Graham (2012), ``How to Get Startup Ideas.''} Emphasizes authentic unmet needs, narrow but growing entry points, founder insight, and tackling unglamorous operational problems.
    \item[\textbf{[R6]}] \textbf{Andrew Chen (2021), \emph{The Cold Start Problem}.} Introduces atomic networks and network density, with direct applications to marketplaces, collaboration tools, and platforms benefiting from network effects.
    \item[\textbf{[R7]}] \textbf{Brian Balfour (2025), ``The Four Fits: A Growth Framework for the AI Era.''} Outlines how product, market, channels, and monetization models must interlock dynamically as a coherent system.
    \item[\textbf{[R8]}] \textbf{Michael E. Porter (1996), ``What Is Strategy?'' \emph{Harvard Business Review}.} Defines strategy as a distinctive, coherent activity system requiring real trade-offs, rather than mere operational efficiency.
    \item[\textbf{[R9]}] \textbf{Chris Zook and James Allen (2011), ``The Great Repeatable Business Model,'' \emph{Harvard Business Review}.} Demonstrates how focusing on a repeatable differentiation model enables sustainable expansion into adjacencies.
    \item[\textbf{[R10]}] \textbf{Y Combinator (2017), ``YC's Essential Startup Advice.''} Practitioner heuristics emphasizing genuine customer pain, manual validation, intense focus, and avoiding premature team expansion.
    \item[\textbf{[R11]}] \textbf{Christensen Institute, ``Disruptive Innovation Theory.''} Contrasts sustaining innovations designed for incumbents' best customers with entrant footholds serving nonconsumers or over-served segments.
    \item[\textbf{[R12]}] \textbf{Previous chapters of \emph{Building What Lasts}.} Chapter 2 provides the founder-asset audit; Chapter 3 defines institutional boundaries; Chapter 4 supplies experiment and option-value reasoning; Chapter 5 establishes market and profit-pool analysis; Chapter 6 provides customer research and buying-behavior methods. This chapter turns those investigations into a deliberately concentrated first commercial move.
\end{enumerate}

\subsection*{The 12 Principles Worth Retaining}

\begin{enumerate}
    \item A beachhead is a \textbf{buyer situation you can win}, not a fashionable niche label.
    \item Narrow by job, trigger, payer, and delivery similarity, not arbitrary demographics.
    \item An underserved segment is one in which alternatives fail a consequential requirement; complaints alone do not prove demand.
    \item A good wedge completes a meaningful job and makes an unknown supplier worth trying.
    \item The best abstract market may be the wrong market for this founder at this stage.
    \item Hard feasibility constraints must be cleared before weighted opportunity comparisons.
    \item A small market can be excellent if the transactions are economically viable and the learning or assets are transferable.
    \item Focus means concentrating \textbf{operations and promises}, not forbidding broad exploration during discovery.
    \item A first sale proves possible demand; repeated independent sales and delivery begin to establish a position.
    \item Expand where proof, customers, channels, processes, or technical assets genuinely transfer, not merely where industries share a label.
    \item Remain committed to the ambition but willing to revise the entry market as evidence arrives.
    \item The initial business should create a position worth building on, rather than a new set of liabilities disguised as growth.
\end{enumerate}

\subsection*{Final Reflection}

The founder who dreams of a large institution may fear that choosing a tiny first market makes the dream smaller. It often does the opposite. It converts the dream into an experiment that can be run, a promise that can be kept, and an asset that can be accumulated.

The beachhead is the first place the market can discover not only what the founder imagines, but \textbf{what the organization can actually do.}

\vspace{1.5ex}
\noindent\emph{End of Chapter 7. Chapter 8 moves from ``Which group do we intend to win first?'' to ``Why should that group choose us, and how should we become the obvious credible alternative?''}
